\section{Algorithmic Checking: \smtlan}
\label{sec:theory:algorithmic}
\label{sec:algorithmic}

% \mypara{Syntax: Terms \& Sorts}
\smtlan is a first order approximation of \corelan
where higher-order features are
approximated with uninterpreted functions and the undecidable
type subsumption rule \rsubbasecore is replaced with a
decidable one (\ie, \rsubbasepbe), yielding an sound
and decidable SMT-based algorithmic type system.
%
Figure~\ref{fig:smtsyntax} summarizes the syntax
of \smtlan, the \emph{sorted} (SMT-)
decidable logic of quantifier-free equality,
uninterpreted functions and linear
arithmetic (QF-EUFLIA) ~\citep{Nelson81,SMTLIB2}.
%
The \emph{terms} of \smtlan include
integers $n$,
booleans $b$,
variables $x$,
data constructors $\dc$ (encoded as constants),
fully applied unary \unop and binary \binop operators,
and application $x\ \overline{\pred}$ of an uninterpreted function $x$.
%
The \emph{sorts} of \smtlan include the built-in
\tint and \tbool to represent
integers and booleans and the 
uninterpreted \tuniv to represent data types.
%
The interpreted functions of \smtlan, \eg
the logical constants $=$ and $<$,
have the function sort $\sort \rightarrow \sort$.
%
Other functional values in \corelan, \eg
reflected \corelan functions and
$\lambda$-expressions, are represented as
first-order values with
the uninterpreted sort \tsmtfun{\sort}{\sort}.


\subsection{Transforming \corelan into \smtlan}
%
\label{subsec:embedding}
\label{subsec:measures}

The judgment
\tologicshort{\env}{e}{\typ}{\pred}{\sort}{\smtenv}{\axioms}
states that a $\corelan$ term $e$ is transformed,
under an environment $\env$, into a
$\smtlan$ term $\pred$.
%
If ${\tologicshort{\env}{\refa}{}{\pred}{}{}{}}$
and $\env$ is clear from the context we write
\tosmt{\refa} and \tohaskell{\pred}
to denote the translation from \corelan to \smtlan and back.
%
Most of the transformation rules are identity
and can be found in~\cite{appendix}.
Here we discuss the non-identity ones.

\mypara{Embedding Types}
%
We embed \corelan types into \smtlan sorts as:
%
$$\arraycolsep=1.4pt
\begin{array}{lclclclclcl}
\tosmt{\tint}   & \defeq &  \tint  & \qquad\qquad &
\tosmt{T}       & \defeq &  \tuniv & \qquad\qquad &
\tosmt{\tref{v}{\btyp^{[\tlabel]}}{e}} &\defeq & \tosmt{\btyp} \\
\tosmt{\tbool}  & \defeq &  \tbool & \quad\quad &
                &        &         & \quad\quad &
\tosmt{\trfun{x}{\typ_x}{\typ}} & \defeq & \tsmtfun{\tosmt{\typ_x}}{\tosmt{\typ}}
\end{array}$$
%
\mypara{Embedding Constants}
%
Elements shared on both \corelan and \smtlan
translate to themselves.
%
These elements include
booleans,
integers,
variables,
binary
and unary
operators.
%
SMT solvers do not support currying,
and so in \smtlan, all function symbols
must be fully applied.
%
Thus, we assume that all applications
to primitive constants and data
constructors are fully applied, \eg by converting
source terms like "(+ 1)" to
"(\z -> z + 1)".
%

\mypara{Embedding Functions}
%
As \smtlan is first-order, we
embed $\lambda$s using the
uninterpreted function \smtlamname{}{}.
%
$$
\inference{
    \tologicshort{\env, \tbind{x}{\typ_x}}{e}{}{\pred}{}{}{} &
        \hastype{\env}{\emptyset}{(\efun{x}{}{e})}{(\trfun{x}{\typ_x}{\typ})}\\
}{
	\tologicshort{\env}{\efun{x}{}{e}}{(\trfun{x}{\typ_x}{\typ})}
	        {\smtlamname{\tosmt{\typ_x}}{\tosmt{\typ}}\ {x}\ {\pred}}
	        {\sort'}{\smtenv, \tbind{f}{\sort'}}{\andaxioms{\{\axioms_{f_1}, \axioms_{f_2}\}}{\axioms}}
}%% [\lgfun]
$$
%
The term $\efun{x}{}{e}$ of type
${\typ_x \rightarrow \typ}$ is transformed
to
${\smtlamname{\sort_x}{\sort}\ x\ \pred}$
of sort
${\tsmtfun{\sort_x}{\sort}}$, where
%
$\sort_x$ and $\sort$ are respectively
$\tosmt{\typ_x}$ and $\tosmt{\typ}$,
%
${\smtlamname{\sort_x}{\sort}}$
is a special uninterpreted function
of sort
${\sort_x \rightarrow \sort\rightarrow\tsmtfun{\sort_x}{\sort}}$,
and
$x$ of sort $\sort_x$ and $\pred$ of sort $\sort$ are
the embeddings of the binder and body, respectively.
%
As $\smtlamname{}{}$ is an SMT-function,
it \emph{does not} create a binding for $x$.
%
Instead, $x$ is renamed to
a \emph{fresh} SMT name.


\mypara{Embedding Applications}
%
We use defunctionalization~\citep{Reynolds72} to 
embed applications. 
% using the uninterpreted $\smtappname{}{}$:
%
$$
\inference{
	\tologicshort{\env}{e_x}{\typ_x}{\pred_x}{\tosmt{\typ_x}}{\smtenv}{\axioms'}
	&
	\tologicshort{\env}{e}{\trfun{x}{\typ_x}{\typ}}{\pred}{\tsmtfun{\tosmt{\typ_x}}{\tosmt{\typ}}}{\smtenv}{\axioms}
	&
	\hastype{\env}{\emptyset}{e}{{\typ_x}\rightarrow{\typ}}
}{
	\tologicshort{\env}{e\ e_x}{\typ}{\smtappname{\tosmt{\typ_x}}{\tosmt{\typ}}\ {\pred}\ {\pred_x}}{\tosmt{\typ}}{\smtenv}{\andaxioms{\axioms}{\axioms'}}
}%% [\lgapp]
$$
%
The term ${e\ e_x}$, where $e$ and $e_x$ have
types ${\typ_x \rightarrow \typ}$ and $\typ_x$,
is transformed to
${\tbind{\smtappname{\sort_x}{\sort}\ \pred\ \pred_x}{\sort}}$
where
%
$\sort$ and $\sort_x$ are $\tosmt{\typ}$ and $\tosmt{\typ_x}$,
the
${\smtappname{\sort_x}{\sort}}$
is a special uninterpreted function of sort
${\tsmtfun{\sort_x}{\sort} \rightarrow \sort_x \rightarrow \sort}$,
and
$\pred$ and $\pred_x$ are the respective translations of $e$ and $e_x$.


\mypara{Embedding Data Types}
%
We embed data constructors to a predefined
\smtlan constant ${\smtvar{\dc}}$ of
sort ${\tosmt{\constty{\dc}}}$:
$
	\tologicshort{\env}{\dc}{\constty{\dc}}{\smtvar{\dc}}{\tosmt{\constty{\dc}}}{\emptyset}{\emptyaxioms}
$.
%
For each datatype, we 
create reflected functions
that \emph{check} the top-level
constructor and \emph{select}
their individual fields.
%
For example, for lists, we
create the functions
%
\begin{mcode}
    isNil []     = True         isCons (x:xs) = True          sel1 (x:xs) = x
    isNil (x:xs) = False        isCons []     = False         sel2 (x:xs) = xs
\end{mcode}
%
The above selectors can be
modeled precisely in the refinement
logic via SMT support for ADTs~\cite{Nelson81}.
%
To generalize, let $\dc_i$ be a
data constructor such that
$
\constty{\dc_i} \defeq \typ_{i,1} \rightarrow \dots \rightarrow \typ_{i,n} \rightarrow \typ
$.
Then \emph{check}
${\checkdc{{\dc_i}}}$ has the sort
$\tsmtfun{\tosmt{\typ}}{\tbool}$
and \emph{select}
${\selector{\dc}{i,j}}$ has the sort
$\tsmtfun{\tosmt{\typ}}{\tosmt{\typ_{i,j}}}$.
%

\mypara{Embedding Case Expressions}
%
We translate case-expressions
of \corelan into nested $\mathtt{if}$
terms in \smtlan, by using the check
functions in the guards and the
select functions for the binders
of each case.
$$
\inference{
	\tologicshort{\env}{e}{\typ_e}{\pred}{\tosmt{\typ_e}}{\smtenv}{\axioms} &&
	\tologicshort{\env}{e_i\subst{\overline{y_i}}{\overline{\selector{\dc_i}{}\ x}}\subst{x}{e}}{\typ}{\pred_i}{\tosmt{\typ}}{\smtenv}{\axioms_i}
}{
	\tologicshort{\env}{\ecase{x}{e}{\dc_i}{\overline{y_i}}{e_i}}{\typ}
	 {\eif{\smtappname{}{}\ \checkdc{\dc_1}\ \pred}{\pred_1}{\ldots} \ \mathtt{else}\ \pred_n}{\tosmt{\typ}}{\smtenv}
	 {\andaxioms{\axioms}{\axioms_i}}
}%[\lgcase]
$$
%
The above translation yields the reflected definition
for append "(++)" from~(\S~\ref{sec:overview:lists}).

\mypara{Semantic Preservation}
%
The translation preserves the semantics of the expressions.
%
Informally,
if ${\tologicshort{\env}{\refa}{}{\pred}{}{}{}}$,
then for every substitution $\theta$ and every logical model $\sigma$
that respects the environment $\env$
if $\evalsto{\applysub{\theta}{\refa}}{v}$
then $\sigma \models \pred = \tosmt{v}$.
%
\VC{What's a logical model?}
%

\subsection{Algorithmic Type Checking}\label{sec:pbe:tc}

We make the type checking from Figure~\ref{fig:typing}
algorithmic by checking subtyping via our novel,
SMT-based \emph{Proof by Logical Evaluation}(\pbesym).
Next, we formalize how \pbesym makes checking algorithmic
and in \S~\ref{sec:pbe} we describe the \pbesym procedure in detail.

\mypara{Verification Conditions}
%
Recall that in ~\S~\ref{subsec:embedding} we defined
$\tosmt{\cdot}$ as the translation from \corelan
to \smtlan.
%
Informally, the implication or \emph{verification condition} (VC)
$\tosmt{\env} \Rightarrow {\pred_1} \Rightarrow {\pred_2}$ is \emph{valid} only if the
set of values described by $\pred_1$ is subsumed by
the set of values described by $\pred_2$ under the
assumptions of \env.
%
$\env$ is embedded into logic by conjoining
the refinements of terminating binders~\cite{Vazou14}:
%
$$
\tosmt{\env} \defeq \bigcup_{x \in \env} \tosmt{\env, x} %\\
\quad \mbox{where we embed each binder as} \quad
\tosmt{\env, x} \defeq \begin{cases}
                           \tosmt{e}  & \text{if } \env(x)=\tref{x}{\btyp^{\tlabel}}{e}\\
                           \etrue & \text{otherwise}.
                         \end{cases}
$$

\mypara{Validity Checking}
Instead of directly using the VCs to check validity of programs, we
use the procedure \pbesym that strengthens the assumption environment \tosmt{\env}
with equational properties.
%
Concretely,
%
given a
%
reflection environment $\RRenv$,
%
type environment $\env$, and
%
expression $e$,
%
the procedure $\pbe{\tosmt{\RRenv}}{\tosmt{\env}}{\tosmt{e}}$
--- we will define $\tosmt{\RRenv}$ in \S~\ref{sec:pbe:algo} ---
returns \ttrue only when the expression $e$ evaluates
to \etrue under the reflection and type environments
\RRenv and \env.

\mypara{Subtyping via VC Validity Checking}
%
We make subtyping, and hence, typing decidable,
by replacing the denotational base subtyping
rule \rsubbasecore with the conservative,
algorithmic version \rsubbasepbe that uses
\pbesym to check the validity of the
subtyping.
%
$$
\inference{
\pbe{\tosmt{\RRenv}}{\tosmt{\env,\tbind{v}{\tref{v}{\btyp^\tlabel}{\refa_1}}}}{\tosmt{\refa_2}}
}{
	\pbeissubtype{\env}{\RRenv}{\tref{v}{\btyp}{\refa_1}}{\tref{v}{\btyp}{\refa_2}}
}[\rsubbasepbe]
$$

This typing rule is sound as functions reflected in \RRenv
always respect the typing environment \env (by construction)
and because \pbename is sound (Theorem~\ref{thm:pbe:sound}).
%
\begin{lemma}\label{lem:subtyping} %[Conservative Subtyping]
If {\pbeissubtype{\env}{\RRenv}{\tref{v}{\btyp}{e_1}}{\tref{v}{\btyp}{e_2}}}
then {\issubtype{\env}{\RRenv}{\tref{v}{\btyp}{e_1}}{\tref{v}{\btyp}{e_2}}}.
\end{lemma}

\mypara{Soundness of \smtlan}
%
We write $\pbehastype{\env}{\RRenv}{e}{\typ}$
for the judgments that can be derived by the
algorithmic subtyping rule \rsubbasepbe (instead of \rsubbasecore.)
%
Lemma~\ref{lem:subtyping} implies the soundness of \smtlan.
%
\begin{theorem}[Soundness of \smtlan]\label{thm:soundness-smt}
If \pbehastype{\env}{\RRenv}{e}{\typ}, then \hastype{\env}{\RRenv}{e}{\typ}.
\end{theorem}
